Wir behaupten

\mavergleichskettealignhandlinks
{\vergleichskettealignhandlinks
{ TY
}
{ =} { \left\{  (P,u)   \mid   h(P) = c , \,  \left(  D h   \right)_{ P } \left(  u  \right) = 0       \right\}
}
{ =} { \left\{  (P,u_1 , \ldots , u_n)   \mid   h(P) = 0 , \,  \sum_{i = 1}^n u_i { \frac{   \partial   h   }{  \partial   x_i   } } (P) = 0       \right\}
}
{ \subseteq} { W \times \R^n
}
{ } {
}
}
{}
{}{}
mit der natürlichen Projektion nach $Y$. Wir fassen dabei $h$ als Funktion auf dem $\R^{2n}$ auf, die nur von den ersten $n$ Variablen abhängt. Die beiden Funktionen $h$  und

\mavergleichskettedisp
{\vergleichskette
{ g(P,u)
}
{ =} { \sum_{i = 1}^n u_i { \frac{   \partial   f   }{  \partial   x_i   } } (P)
}
{ } {
}
{ } {
}
{ } {
}
}
{}{}{}
sind stetig und daher ist eine durch Gleichungen bestimmte abgeschlossene Teilmenge von $W \times \R^n$ gegeben. Zu jedem festen Punkt

\mavergleichskette
{\vergleichskette
{ P
}
{ \in }{  Y
}
{  }{
}
{    }{
}
{   }{
}
}
{}{}{}
ist

\mavergleichskettedisp
{\vergleichskette
{  \pi^{-1}(P)
}
{ =} { \left\{  (u_1 , \ldots , u_n)   \mid   \sum_{i = 1}^n u_i { \frac{   \partial   h   }{  \partial   x_i   } } (P) = 0         \right\}
}
{ \subseteq} {  \R^n
}
{ } {
}
{ } {
}
}
{}{}{}
der
\definitionsverweis {Tangentialraum an die Faser}{}{}
$Y$ im Punkt $P$. Aufgrund es Satzes über implizite Abbildungen gibt es zu jedem Punkt

\mavergleichskette
{\vergleichskette
{ P
}
{ \in }{  Y
}
{  }{
}
{    }{
}
{   }{
}
}
{}{}{}
eine offene Umgebung

\mavergleichskette
{\vergleichskette
{ P
}
{ \in }{  U
}
{ \subseteq }{ Y
}
{    }{
}
{   }{
}
}
{}{}{,}
eine offene Teilmenge

\mavergleichskette
{\vergleichskette
{ V
}
{ \subseteq }{  \R^{n-1}
}
{  }{
}
{    }{
}
{   }{
}
}
{}{}{}
und ein Homöomorphismus
\maabbdisp {\beta} { V } {U
} {,}
der als Abbildung nach $\R^n$ stetig differenzierbar ist und eine stetige in der zweiten Komponenten lineare Abbildung
\maabbeledisp {} {V \times \R^{n-1} } { U \times \R^n
} {(Q,z)} { ( \beta(q), \left(  D\beta  \right)_{ Q } \left(  z  \right) )
} {,}
induziert, die stets im Tangentialraum an die Faser landet. Dies zeigt, dass die eingebettete Realisierung auch mit der Bündelstruktur und der Topologie des Tangentialbündels übereinstimmt.

 [[Kategorie:Latexseite]]
